\chapter{LASIK surgery}

\section*{Definition and Overview}
Laser-Assisted in Situ Keratomileusis (LASIK) is a highly effective refractive surgical procedure designed to correct common vision problems, including myopia (nearsightedness), hyperopia (farsightedness), and astigmatism. The procedure involves using an advanced laser (femtosecond laser) or a mechanical microkeratome to create a thin corneal flap, which is then lifted to allow an excimer laser to precisely reshape the underlying corneal stroma. By reshaping the cornea, LASIK alters its refractive power, enabling light entering the eye to focus properly onto the retina without the aid of glasses or contact lenses. As a minimally invasive, outpatient ophthalmic procedure, LASIK is widely celebrated for its rapid visual recovery, minimal postoperative discomfort, and high rates of patient satisfaction. It has revolutionised the field of refractive surgery, offering millions of individuals globally a path toward functional visual independence.

\section*{Epidemiology}
LASIK is one of the most frequently performed elective procedures worldwide, with millions of surgeries completed since its approval in the late 1990s. In Australia and other developed countries, hundreds of thousands of refractive procedures are performed annually, with LASIK comprising the vast majority of these cases. The target demographic typically ranges from 18 to 45 years of age, matching the period of refractive stability before the onset of age-related presbyopia. There is no significant gender predisposition, with males and females seeking the procedure in comparable proportions. While hyperopia and astigmatism corrections are common, myopia correction represents the largest cohort of LASIK patients, reflecting the rising global prevalence of nearsightedness. Postoperative satisfaction rates exceed 95\%, while major long-term visual complications occur in fewer than 1\% of carefully screened patients.

\section*{Aetiology and Pathophysiology}
The physiological need for LASIK arises from refractive errors, which occur when the eye's shape prevents light from focusing directly on the retina.
\begin{itemize}
    \item \textbf{Myopia (nearsightedness):} A refractive state where the cornea is too steeply curved or the eyeball is too long, causing light rays to focus in front of the retina instead of directly on it.
    \item \textbf{Hyperopia (farsightedness):} A refractive state where the cornea is too flat or the eyeball is too short, causing light rays to focus behind the retina.
    \item \textbf{Astigmatism:} An irregular curvature of the cornea or lens, resembling a rugby ball rather than a basketball, leading to multiple focal points and blurred vision at all distances.
    \item \textbf{Corneal stromal remodelling:} The excimer laser corrects these errors by photoablation, breaking carbon-carbon molecular bonds to remove microscopic layers of corneal tissue without generating heat.
    \item \textbf{Biomechanical response:} Creating a corneal flap weakens the anterior corneal lamellae, necessitating strict preoperative screening to ensure the residual stromal bed is thick enough to prevent corneal ectasia.
\end{itemize}

\section*{Clinical Features}
Patients seeking LASIK present with symptoms of refractive errors, while the postoperative phase is characterised by specific recovery features.
\begin{itemize}
    \item \textbf{Refractive blurring:} Difficulty seeing distant objects clearly (myopia), close-up objects clearly (hyperopia), or distortion at all distances (astigmatism).
    \item \textbf{Asthenopia:} Eyestrain, headaches, and physical fatigue associated with chronic squinting or focusing efforts.
    \item \textbf{Immediate postoperative burning:} A mild, transient gritty or burning sensation in the eyes for 3--4 hours after the anaesthetic drops wear off.
    \item \textbf{Fluctuating vision:} Blurred or unstable vision during the first few days to weeks as the corneal flap adheres and epithelial healing progresses.
    \item \textbf{Halos and glare:} Temporary visual disturbances, such as starbursts, glare, and halos around lights at night, which typically resolve over 1--3 months.
\end{itemize}

\section*{Differential Diagnosis}
In the preoperative phase, other causes of visual impairment must be distinguished from simple refractive errors. Postoperatively, complications must be differentiated from normal healing.
\begin{itemize}
    \item \textbf{Keratoconus:} A progressive thinning and cone-like protrusion of the cornea, which is an absolute contraindication for LASIK.
    \item \textbf{Presbyopia:} Age-related loss of lens accommodation, which is distinct from corneal refractive errors and not corrected by standard bilateral distance-focused LASIK.
    \item \textbf{Diffuse lamellar keratitis (DLK):} A sterile inflammatory reaction under the corneal flap, characterised by a "sands of Sahara" appearance on slit-lamp exam, which must be distinguished from microbial keratitis.
    \item \textbf{Microbial keratitis:} A rare but serious corneal infection requiring aggressive antibiotic therapy rather than steroid adjustments.
    \item \textbf{Dry eye syndrome:} Postoperative tear film instability due to corneal nerve transection, which must be differentiated from pre-existing chronic dry eyes.
\end{itemize}

\section*{When to Seek Medical Care}
Patients should seek a consultation with an ophthalmologist to determine eligibility. Following surgery, prompt evaluation is required if any signs of infection or flap displacement occur.

\textbf{Call 000 (or local emergency services) for:}
\begin{itemize}
    \item \textbf{Anaphylactic reactions:} Severe allergic reaction to prescribed eye drops (e.g. difficulty breathing, swelling of the face or throat).
    \item \textbf{Severe eye trauma:} Direct physical impact to the eye post-surgery that could dislodge the corneal flap.
\end{itemize}

Other reasons to contact your ophthalmologist urgently include:
\begin{itemize}
    \item \textbf{Sudden decrease in vision:} A significant drop in visual acuity or progressive blurring.
    \item \textbf{Severe or worsening eye pain:} Discomfort that is not controlled by standard over-the-counter pain relief.
    \item \textbf{Purulent discharge:} Thick, yellow or green discharge from the eye, indicating infection.
\end{itemize}

\section*{Diagnosis and Investigations}
A rigorous preoperative assessment is vital to ensure patient safety and optimal visual outcomes.
\begin{itemize}
    \item \textbf{Corneal topography and tomography:} Advanced mapping of the anterior and posterior corneal surfaces to assess curvature, thickness, and rule out subclinical keratoconus.
    \item \textbf{Pachymetry:} Ultrasonic or optical measurement of corneal thickness to ensure sufficient residual stroma remains after photoablation.
    \item \textbf{Dry eye evaluation:} Assessment of tear film stability and volume using Schirmer's tests or tear breakup time (TBUT) measurements.
    \item \textbf{Manifest and cycloplegic refraction:} Precise measurements of refractive error before and after pharmacological relaxation of the ciliary muscle.
    \item \textbf{Pupillometry:} Measurement of pupil size in dark and light conditions to assess the risk of postoperative night vision disturbances.
\end{itemize}

\section*{Management}
The management of a LASIK patient includes precise surgical execution and structured postoperative care.
\begin{itemize}
    \item \textbf{Surgical flap creation:} Creating a corneal flap using a femtosecond laser, which uses ultra-short pulses of light to separate tissue layers at a pre-programmed depth.
    \item \textbf{Excimer laser ablation:} Applying computer-guided ultraviolet laser pulses to vaporise stromal tissue and reshape the cornea, guided by preoperative wavefront maps.
    \item \textbf{Pharmacological eye drops:} Postoperative administration of prophylactic antibiotic drops (e.g. moxifloxacin) and anti-inflammatory corticosteroid drops (e.g. prednisolone acetate).
    \item \textbf{Ocular lubrication:} Frequent use of preservative-free artificial tears to support corneal epithelial healing and manage transient dry eyes.
    \item \textbf{Protective eyewear:} Wearing clear plastic shields during sleep for the first week to prevent accidental eye rubbing.
\end{itemize}

\section*{Complications}
\begin{itemize}
    \item \textbf{Dry eye syndrome:} Transient or chronic decrease in tear production due to temporary damage to the corneal nerves during flap creation.
    \item \textbf{Flap complications:} Partial flaps, free caps, striae (wrinkles in the flap), or flap displacement, which may require surgical repositioning.
    \item \textbf{Under-correction or over-correction:} Inaccurate laser ablation resulting in residual refractive error, which may require a secondary enhancement procedure.
    \item \textbf{Epithelial ingrowth:} Migration of corneal epithelial cells under the flap during healing, which can cause astigmatism or visual distortion.
    \item \textbf{Corneal ectasia:} Progressive thinning and bulging of the cornea due to biomechanical instability, mimicking keratoconus.
\end{itemize}

\section*{Prognosis and Outcomes}
The prognosis for patients undergoing LASIK is highly favourable. The vast majority of patients (over 95\%) achieve 20/20 uncorrected visual acuity or better, and more than 99\% achieve at least 20/40, which is the standard legal limit for driving without corrective lenses. Visual recovery is remarkably rapid, with many patients noticing a significant improvement within 24 hours of surgery. Mild visual side effects, such as dry eyes and night-time glare, are common in the early postoperative weeks but typically resolve or stabilise within 3 to 6 months. While LASIK does not prevent the natural progression of presbyopia (the age-related need for reading glasses) or cataracts later in life, it provides a stable, long-lasting reduction in dependence on glasses or contact lenses for distance vision.

\section*{Prevention and Self-Care}
\begin{itemize}
    \item \textbf{Avoid rubbing the eyes:} Refraining from rubbing the eyes for at least 1 month post-surgery to prevent corneal flap displacement.
    \item \textbf{Use prescribed drops diligently:} Following the exact schedule for antibiotic, steroid, and lubricating drops.
    \item \textbf{Wear eye protection:} Using sunglasses with UV protection outdoors and safety glasses during contact sports or dusty activities.
    \item \textbf{Avoid water exposure:} Staying out of swimming pools, hot tubs, and natural bodies of water for at least 2 weeks to reduce infection risks.
    \item \textbf{Limit screen time initially:} Resting the eyes for the first 24--48 hours by limiting television, computer, and phone use to prevent dryness.
\end{itemize}

\section*{Sources and Further Reading}
The following reputable organisations provide further information on LASIK surgery and refractive care guidelines:
\begin{itemize}
    \item American Academy of Ophthalmology. \textit{Patient guides and clinical updates on LASIK surgery.} https://www.aao.org/
    \item Royal Australian and New Zealand College of Ophthalmologists. \textit{Patient resources on refractive surgery in Australia.} https://ranzco.edu/
    \item Mayo Clinic. \textit{What to expect during and after LASIK surgery.} https://www.mayoclinic.org/
    \item U.S. Food and Drug Administration. \textit{LASIK information, risks, and checklist for patients.} https://www.fda.gov/
    \item National Health Service. \textit{Overview of laser eye surgery and refractive options.} https://www.nhs.uk/
\end{itemize}
